%% kapa: exam=2 round=27 year=2016 period=3 subject=compensation_law subject_ko=감정평가 및 보상법규 seq=3 number=3 points=20
%% source: Q-Net 감정평가사 기출문제 (cst00309) artlSeq=5185401 file=extracted/1008/5185401/제27회 감정평가사 2차시험 기출문제/제27회 감정평가사 2차 시험문제지/(3교시 시험문제지)감정평가 및 보상법규.pdf
%% conversion: kordoc 4.24.51 -> tools/gen2.py (auto); review=proofread-gpt-sol; flags=none
\begin{longquestion}{【문제 3】}
국방부장관은 국방ㆍ군사에 관한 사업을 위하여 국토교통부장관으로부터 甲 소유의 토지를 포함한 200필지의 토지 600,000m\textsuperscript{2}에 관하여 「공익사업을 위한 토지 등의 취득 및 보상에 관한 법률」제20조에 따른 사업인정을 받았다. 그러나 국토교통부장관은 사업인정을 하면서 동법 제21조에 규정된 이해관계인의 의견을 청취하는 절차를 거치지 않았다. 한편, 국방부장관은 甲과 손실보상 등에 관하여 협의하였으나 협의가 성립되지 않았다. 국방부장관은 재결을 신청하였고 중앙토지수용위원회는 수용재결을 하였다. 甲은 수용재결에 대한 취소소송에서 사업인정의 절차상 하자를 이유로 수용재결의 위법성을 주장할 수 있는가? (단, 국토교통부장관의 사업인정에 대한 취소소송의 제소기간은 도과하였음) (20점) \par
\end{longquestion}
